
\chapter{The Observed-State Safety Illusion}
\label{ch:safety-illusion}

A battery controller that checks constraints only against its own
observations will, by construction, never detect a violation it cannot
see.  This chapter formalises the \emph{observed-state safety illusion},
shows how it arises in the ORIUS dispatch loop, explains why neither
better forecasting nor tighter optimisation can resolve it, and introduces
the dual-trajectory evaluation paradigm that makes the illusion visible.

\section{From Partial Observability to Safety Illusion}

Partial observability is a well-studied concept in control theory and
reinforcement learning.  A partially observable Markov decision process
(POMDP) acknowledges that the agent sees a noisy or incomplete projection
of the true state and plans accordingly, typically via belief-state
updates.  The safety illusion arises when the agent \emph{does not plan
accordingly}---when the controller treats its observation as ground truth
and evaluates constraints against it without accounting for observation
quality.

In a POMDP with a well-calibrated belief, the agent maintains a
distribution over true states and can compute the worst-case constraint
satisfaction probability.  In practice, battery dispatch controllers do
not maintain belief distributions; they receive a point estimate
$\hat{x}_t$ from the SCADA system and optimise against it.  The gap
between the POMDP ideal and the deployed reality is the source of the
illusion.

\section{Why the Illusion Arises in the ORIUS Loop}

\subsection{The Optimise-to-Dispatch Boundary}

As described in Chapter~\ref{ch:problem-formulation}, the ORIUS loop
has a clear boundary between the optimiser and the dispatcher.  The
optimiser solves~\eqref{eq:dispatch-obj} on the observed state and
produces a candidate action~$a_t^{\star}$.  The dispatcher (in the
baseline case) forwards this action to the plant without modification.

Both the optimiser and the baseline dispatcher operate on
$\hat{x}_t$.  The safety constraints
\[
  \mathrm{SOC}_{\min} \;\le\;
  \widehat{\mathrm{SOC}}_t + \frac{\eta_c [a_t]^{+} - [a_t]^{-}/\eta_d}{E_{\max}} \cdot \Delta t
  \;\le\; \mathrm{SOC}_{\max}
\]
are checked on the observed SOC.  If $\widehat{\mathrm{SOC}}_t$ over- or
under-estimates the true SOC, the constraint check is meaningless with
respect to physical safety.

\subsection{Why Good Forecasting Does Not Solve It}

A natural objection is: ``Improve the forecast, and the observation error
disappears.''  This is incorrect for a fundamental reason.  The forecast
model predicts future load, solar, and wind \emph{given current
observations}.  If the current observation is itself degraded---a stale
SOC reading, a dropped load packet---the forecast inherits the error.  A
perfect forecast model with imperfect inputs produces imperfect outputs.

More precisely, the forecast model learns $\mathbb{E}[Y_{t+h} \mid
\hat{X}_t]$, not $\mathbb{E}[Y_{t+h} \mid X_t]$.  When
$\hat{X}_t \ne X_t$, the conditional expectation is evaluated at the wrong
conditioning point.  The resulting action is optimal for a state the plant
does not occupy.

\section{Formal Statement of the Illusion}
\label{sec:formal-illusion}

\begin{definition}[Observed-State Safety Illusion]
\label{def:safety-illusion}
Let $\hat{x}_t$ be the observed state, $x_t$ the true state, and
$a_t$ the dispatched action.  The \emph{observed-state safety illusion}
at time~$t$ is the conjunction of three conditions:
\begin{enumerate}
  \item \textbf{Observed safety:}
    $a_t \in \mathcal{A}(\hat{x}_t)$---the action satisfies all
    constraints when evaluated on the observation.
  \item \textbf{True-state violation:}
    $a_t \notin \mathcal{A}(x_t)$---the action violates at least one
    constraint when evaluated on the truth.
  \item \textbf{Undetectability:}
    The controller has no mechanism to detect condition~(2) from its
    available information $\hat{x}_t$ alone.
\end{enumerate}
\end{definition}

Condition~(3) distinguishes the illusion from a detectable fault.  If the
controller could infer that $\hat{x}_t$ is unreliable (e.g., via an OQE
score), it could compensate.  The illusion is \emph{silent}: the
controller proceeds with full confidence in a physically unsafe action.

\begin{theorem}[Existence of the illusion under dropout]
\label{thm:illusion-existence}
Under the fault model of Section~\ref{sec:fault-model} with dropout
rate $\rho > 0$ and last-value hold imputation, there exists a time
step~$t$ and a deterministic LP dispatch action~$a_t$ such that the
observed-state safety illusion (Definition~\ref{def:safety-illusion})
holds with positive probability.
\end{theorem}

\begin{proof}[Proof sketch]
Consider a discharge phase where
$\mathrm{SOC}_t^{\mathrm{true}} = \mathrm{SOC}_{\min} + \epsilon$ for
small $\epsilon > 0$.  If the most recent $k$ SOC readings were dropped
and the last-value hold reports
$\widehat{\mathrm{SOC}}_t = \mathrm{SOC}_{\min} + \epsilon + k \cdot
\delta$ (the SOC at the time of the last successful reading, before $k$
discharge steps each consuming $\delta$), then the LP observes a
comfortable margin and dispatches a further discharge~$a_t < 0$.  The
true SOC after this action is
$\mathrm{SOC}_{t+1}^{\mathrm{true}} = \mathrm{SOC}_{\min} + \epsilon -
|a_t| \cdot \Delta t / (E_{\max} \cdot \eta_d) < \mathrm{SOC}_{\min}$
for $|a_t|$ exceeding $\epsilon \cdot E_{\max} \cdot \eta_d / \Delta t$.
The action satisfies observed constraints, violates true constraints, and
the LP cannot detect the violation---all three conditions hold.
\end{proof}

\section{Battery-Domain Realisation Under Telemetry Faults}
\label{sec:battery-narrative}

\subsection{A Simple Battery Narrative}

Consider a 100~MWh / 50~MW battery operating over 168~hours (one week)
under the \texttt{dropout\_only} scenario at $\rho = 0.20$.

\begin{enumerate}
  \item \textbf{Hour 0--24.}  Telemetry is clean.  The deterministic LP
    dispatches charge during solar surplus and discharge during evening
    peak.  Observed and true SOC tracks are indistinguishable.
  \item \textbf{Hour 25.}  A dropout burst begins: five consecutive SOC
    readings are missing.  The last-value hold freezes
    $\widehat{\mathrm{SOC}}$ at 55\%.
  \item \textbf{Hour 26--28.}  The LP continues discharging against the
    frozen observation.  True SOC falls from 55\% to 38\% to 24\%.  The
    observed SOC remains at 55\%.
  \item \textbf{Hour 29.}  The LP dispatches a 40~MW discharge, believing
    SOC is 55\%.  True SOC drops to 16\%, breaching the 20\% floor.
    Observed SOC still shows 55\%.  The illusion is complete.
  \item \textbf{Hour 30.}  Telemetry recovers.  The LP observes SOC at
    16\% and abruptly curtails dispatch, but the damage---an SOC
    excursion below the safety floor---has already occurred.
\end{enumerate}

This narrative is not hypothetical.  The CPSBench harness reproduces it
deterministically (seed 42, \texttt{dropout\_only}, $\rho = 0.20$), and
the true-state violation is recorded in the TSVR metric.

\section{Why Standard Evaluation Misses the Problem}
\label{sec:standard-eval-fails}

Standard battery dispatch evaluation protocols report:

\begin{itemize}
  \item Revenue or regret (economic performance).
  \item Observed SOC constraint violations (observed safety).
  \item Forecast accuracy (MAPE, RMSE, CRPS).
\end{itemize}

None of these metrics involve the \emph{true} SOC trajectory.  A
controller that achieves zero observed violations can simultaneously have
a high true-state violation rate.  The evaluation is complicit in the
illusion: it checks safety on the same degraded observation that caused
the unsafe action.

\begin{proposition}[Insufficiency of observed-state evaluation]
\label{prop:obs-eval-insufficient}
For any $\mathrm{TSVR} > 0$, there exists a controller-scenario pair
with $\mathrm{OSVR} = 0$ and the given $\mathrm{TSVR}$.  Observed-state
evaluation is therefore \emph{insufficient} to certify physical safety.
\end{proposition}

\begin{proof}
Take any controller $\pi$ and any fault scenario with dropout rate
$\rho > 0$.  Since $\pi$ evaluates constraints on $\hat{x}_t$ and the
telemetry operator preserves observed-state feasibility by construction
(the hold value is always within bounds), $\mathrm{OSVR} = 0$.  By
Theorem~\ref{thm:illusion-existence}, there exists a step where $a_t$
violates true-state constraints, so $\mathrm{TSVR} > 0$.
\end{proof}

\section{CPSBench and Dual-Trajectory Evaluation}
\label{sec:cpsbench-dual}

To break the illusion, we introduce the \emph{dual-trajectory} evaluation
paradigm implemented in the CPSBench harness.  At each step~$t$, the
harness maintains two parallel state machines:

\begin{enumerate}
  \item \textbf{Observed trajectory.}  The SOC evolution computed from
    the action~$a_t$ and the observed state~$\hat{x}_t$, using the
    dynamics~\eqref{eq:soc-dynamics} with $\widehat{\mathrm{SOC}}_t$ as
    the initial condition.
  \item \textbf{True trajectory.}  The SOC evolution computed from the
    \emph{same} action~$a_t$ but with the true
    state~$\mathrm{SOC}_t^{\mathrm{true}}$ as the initial condition.
\end{enumerate}

The true trajectory is never visible to the controller; it is maintained
solely for evaluation.  This separation is critical: the controller
cannot ``cheat'' by observing the truth, and the evaluator cannot be
fooled by the observation.

\begin{definition}[True-State Violation Rate (TSVR)]
\label{def:tsvr}
\[
  \mathrm{TSVR} \;=\; \frac{1}{T} \sum_{t=1}^{T}
  \mathbf{1}\!\Bigl[\,
    \mathrm{SOC}_{t}^{\mathrm{true}} < \mathrm{SOC}_{\min}
    \;\;\text{or}\;\;
    \mathrm{SOC}_{t}^{\mathrm{true}} > \mathrm{SOC}_{\max}
  \,\Bigr].
\]
\end{definition}

\begin{definition}[Intervention Rate (IR)]
\label{def:ir}
\[
  \mathrm{IR} \;=\; \frac{1}{T} \sum_{t=1}^{T}
  \mathbf{1}\!\bigl[\,a_t^{\mathrm{safe}} \ne a_t^{\star}\,\bigr],
\]
where $a_t^{\star}$ is the optimiser's proposed action and
$a_t^{\mathrm{safe}}$ is the action actually dispatched after the
DC3S shield.
\end{definition}

The dual-trajectory architecture enables unambiguous attribution: every
true-state violation is traceable to a specific action, a specific
observation error, and a specific gap between $\hat{x}_t$ and $x_t$.

\section{Why the Illusion Demands a Runtime Shield}
\label{sec:shield-motivation}

The illusion cannot be resolved at the optimiser level for two reasons:

\begin{enumerate}
  \item \textbf{Information asymmetry.}  The optimiser has access only to
    $\hat{x}_t$.  Without a reliability signal, it cannot distinguish
    clean telemetry from degraded telemetry.  Any fixed margin it
    applies will be either too conservative (wasting revenue when
    telemetry is clean) or too aggressive (missing violations when
    telemetry degrades).
  \item \textbf{Separation of concerns.}  The optimiser's objective is
    economic; the safety layer's objective is physical.  Mixing the two
    creates Pareto conflicts that are better resolved by a dedicated
    post-optimiser shield.  The shield can adapt its tightening in
    real time based on the OQE score, something a batch optimiser
    solved minutes ago cannot do.
\end{enumerate}

The DC3S shield resolves the illusion by introducing an
\emph{observation-quality-aware} constraint tightening that operates at
the dispatch boundary.  The shield does not replace the optimiser; it
corrects the optimiser's blind spot.

\begin{theorem}[DC3S Feasibility Guarantee]
\label{thm:dc3s-informal}
Let $\alpha \in (0,1)$ be the conformal miscoverage level, $n$ the
calibration-set size, and $\bar{w} = T^{-1}\sum_{t=1}^T w_t$ the
episode-mean reliability score.  Under Assumptions~A1, A2, A3, and~A5,
the DC3S shield with conformal quantile $\hat{q}_\alpha$ inflated by
$1/w_t$ satisfies, for any finite episode of $T$ steps:
\[
  \mathrm{TSVR} \;\le\; \alpha\,(1 - \bar{w}) \;+\; \varepsilon_n,
\]
where $\varepsilon_n = O(1/\sqrt{n})$ is the finite-sample calibration
correction term.  In particular:
\begin{enumerate}
  \item At $\bar{w} = 1$ (perfect telemetry) and large $n$:
    $\mathrm{TSVR} \to 0$.
  \item At $\bar{w} < 1$ (degraded telemetry): the bound is
    proportional to the mean reliability gap $(1 - \bar{w})$.
  \item As $n \to \infty$: $\varepsilon_n \to 0$ and the bound
    converges to the episode-mean reliability term exactly.
\end{enumerate}
This result is the single-episode analogue of Theorem~T3
(ORIUS Core Bound, Chapter~\ref{ch:core-bound}), which proves
$\mathbb{E}[V] \le \alpha(1-\bar{w})T$ in expectation over episodes.
\end{theorem}

The formal version of this theorem, with precise assumptions and a
complete proof, is developed in subsequent chapters.

\section{Chapter Summary}

The observed-state safety illusion is the mechanism by which the OASG
operates in practice: a controller evaluates constraints on a degraded
observation, concludes the system is safe, and dispatches an action that
is physically unsafe.  The illusion is invisible to any evaluation
protocol that checks constraints on the observed trajectory only.  The
dual-trajectory architecture of CPSBench makes the illusion visible, and
the DC3S runtime shield makes it preventable.  The remainder of the
dissertation develops the shield, proves its guarantees, and validates
them empirically across US and German grid data under the full fault model
of Chapter~\ref{ch:problem-formulation}.
